\section{A regular character chart for the same piece}\label{ng:a-regular-character-chart-for-the-same-piece}

The face chart identifies the actual topology. The following rational
matrix chart is useful both for its peripheral words and for continuation
to generic auxiliary values over any fixed \(M\ne0\).
Throughout this section use

\begin{equation}\label{ng:C0}
\begin{gathered}
 t=M+M^{-1},\quad D=M^2+M^{-2}=t^2-2,\\
 \beta=Q+Q^{-1},\quad\gamma=P+P^{-1},\\
 A_0=X+X^{-1},\quad S=\beta+\gamma+t^2.
 \end{gathered}
\end{equation}

Equation \eqref{ng:T1} is

\begin{equation}\label{ng:C0a}
(A_0+\beta)(A_0+\gamma)=(A_0-2)(A_0-D),\qquad
 A_0=\frac{2D-\beta\gamma}{S}.
\end{equation}

The denominator is eight at completeness. In the following matrix
formulas the internal mixed trace is denoted \(\mathcal X\), so it
is distinct from the source wire \(X\).

\begin{theorem}[The local character and its marking]\label{ng:character}
The matrices constructed in this section give a determinant-one
four-puncture representation and invertible projective auxiliary
transports on their displayed regular open set of \eqref{ng:T1}.
Their full two-generator conjugacy identities, primitive auxiliary
ratios and local main commutator are rational identities there.
Near the complete point they are conjugate, with the indicated
selected flags and original peripheral words, to the actual face
representation of Section~\ref{ng:the-local-face-maps}.
\end{theorem}

We prove the statement by constructing the matrices and then matching
all their words to the native cut.

\subsection{A four-puncture character and its explicit matrices}\label{ng:a-four-puncture-character-and-its-explicit-matrices}

Put

\begin{equation}\label{ng:C5}
p=-\beta,\qquad q=-\gamma,\qquad \mathcal X=2+SX,
 \qquad b=p-2,\quad h=M-M^{-1}.
\end{equation}

For the moment \(p,q,\mathcal X\) may be arbitrary nearby parameters. Define

\begin{equation}\label{ng:C6}
A=\begin{pmatrix}M&1\\0&M^{-1}\end{pmatrix},
 \qquad
 B=\begin{pmatrix}M^{-1}&0\\b&M\end{pmatrix}.
\end{equation}

Both have determinant one and trace \(t\), and \(\operatorname{tr}(AB)=p\). Set

\begin{equation}\label{ng:C7}
\begin{gathered}
 c=\frac{M\mathcal X+M^{-1}q-2t}{D-p},\qquad
 d=\frac{q-M^2-1+hc}{p-2},\\
 e=\mathcal X-1-M^{-2}-hc,\qquad
 C=\begin{pmatrix}c&d\\e&t-c\end{pmatrix}.
 \end{gathered}
\end{equation}

The denominators \(p-2=-4\) and \(D-p=4\) at the complete point are nonzero. The construction does not divide by \(M-M^{-1}\). Direct multiplication gives

\begin{equation}\label{ng:C8}
\operatorname{tr}C=t,\qquad
 \operatorname{tr}(BC)=q,\qquad
 \operatorname{tr}(AC)=\mathcal X,\qquad
 \operatorname{tr}(ABC)=t.
\end{equation}

For the last identity, the nonconstant part of the trace is
\(bc+Me+M^{-1}bd\), which vanishes by \eqref{ng:C7}.

The determinant calculation is

\begin{equation}\label{ng:C9}
\det C-1=\frac{\mathcal F(\mathcal X,p,q)}{(p-2)(D-p)},
\end{equation}

where

\begin{equation}\label{ng:C10}
\mathcal F=
 \mathcal X^2+p^2+q^2+pq\mathcal X-2t^2(p+q+\mathcal X)+4t^2+t^4-4.
\end{equation}

This is the equal-boundary-trace four-punctured-sphere character equation, derived here from the matrices. For a direct verification of \eqref{ng:C9}, substitute the three entries in \eqref{ng:C7}, multiply by \((p-2)(D-p)\), and collect the quadratic polynomial in \(\mathcal X\). Its leading coefficient is one; the other two coefficients are those displayed in \eqref{ng:C10}.

Now impose \eqref{ng:C5}. A direct substitution gives the especially simple identity

\begin{equation}\label{ng:C11}
\mathcal F(2+SX,-\beta,-\gamma)
 =XS\bigl[(A_0+\beta)(A_0+\gamma)
                         -(A_0-2)(A_0-D)\bigr].
\end{equation}

Thus \eqref{ng:T1} is equivalent to \(\det C=1\) on this regular chart. Put

\begin{equation}\label{ng:C12}
D_0=(ABC)^{-1}.
\end{equation}

Equations \eqref{ng:C8}--\eqref{ng:C12} give

\begin{equation}\label{ng:C13}
ABCD_0=I,\qquad
 \det A=\det B=\det C=\det D_0=1,\qquad
 \operatorname{tr}A=\operatorname{tr}B
       =\operatorname{tr}C=\operatorname{tr}D_0=t.
\end{equation}

These are literal matrices representing the four positively oriented boundary generators of a four-punctured sphere. Each puncture has eigenvalues \(M,M^{-1}\). In particular, the source relation is not merely similar in degree to a character equation: \eqref{ng:C7}, \eqref{ng:C11}, and \eqref{ng:C12} realize it by actual matrices.

At the complete selected point,

\begin{equation}\label{ng:C14}
A=\begin{pmatrix}1&1\\0&1\end{pmatrix},\quad
 B=\begin{pmatrix}1&0\\-4&1\end{pmatrix},\quad
 C=\begin{pmatrix}-1+2i&1\\8i&3-2i\end{pmatrix}.
\end{equation}

All four punctures are noncentral parabolics. The pair \(A,B\) has no common eigenline. The construction is therefore on a regular irreducible character chart through the complete point.

The alternative root of the quadratic \eqref{ng:C10} is \(2+S/X\). Indeed its sum with \(\mathcal X=2+SX\) is \(2t^2-\beta\gamma\), by \eqref{ng:C0a}. For comparison, its discriminant factors as

\begin{equation}\label{ng:C15}
(\beta+2)(\gamma+2)
       \bigl[\beta\gamma-2\beta-2\gamma-4D-4\bigr]
           =S^2(X-X^{-1})^2.
\end{equation}

This verifies the two-sheeted parameter and its branch without assigning signs by a numerical square root.

\subsection{Auxiliary transport matrices with the actual source ratios}\label{ng:auxiliary-transport-matrices-with-the-actual-source-ratios}

Define two invertible projective transport matrices

\begin{equation}\label{ng:C16}
V=AB+z_v I,\qquad T=BC+z_w I,
 \qquad
 z_v=-\frac{1+X(\beta+1)}{1-X},\qquad
 z_w=\frac{X+\gamma+1}{1-X}.
\end{equation}

Scalar multiples of \(V,T\) give the same projective transports. Their determinant-one normalizations can be chosen on a local square-root cover; every commutator below is independent of that choice.

The following full matrix identities hold on \eqref{ng:T1}:

\begin{equation}\label{ng:C17}
\begin{aligned}
 VA&=D_0^{-1}V,& VB&=C^{-1}V,\\
 TB&=A^{-1}T,& TC&=D_0^{-1}T.
 \end{aligned}
\end{equation}

They are stronger than trace matching. Also \(V\) commutes with \(AB\), and \(T\) commutes with \(BC\), directly from \eqref{ng:C16}.

Here is a rational way to verify \eqref{ng:C17}, preserving the regular chart. Eliminate \(\gamma\) using \eqref{ng:C0a}:

\begin{equation}\label{ng:C18}
\gamma=\frac{2D-(X+X^{-1})(\beta+D+2)}
                         {X+X^{-1}+\beta}.
\end{equation}

The denominator is two at the complete point. Substitute \eqref{ng:C18}, \eqref{ng:C5}, and \eqref{ng:C7} in \eqref{ng:C17}. Every difference becomes a zero rational function of \(M,\beta,X\). This is also a polynomial certificate after clearing the recorded denominators. For example, before substituting \eqref{ng:C5}, the value of \(z\) solving the lower-left entry of
\((AB+zI)A=(ABC)(AB+zI)\) makes the remaining three entries scalar multiples of \eqref{ng:C10}; after \eqref{ng:C5}, that \(z\) reduces to \(z_v\). The other pair follows by cyclic permutation of the four punctures. No eigenline diagonalization at a repeated eigenvalue is used.

The two disk-longitude words and their selected eigenvalues are

\begin{equation}\label{ng:C19}
\Lambda_v=AB,\quad \lambda_v=-Q;
 \qquad
 \Lambda_w=(BC)^{-1},\quad \lambda_w=-P.
\end{equation}

The inverse in the second word records the opposite boundary orientation. On the \(-Q\) eigenline of \(AB\), the projective eigenvalue ratio of \(V\) is

\begin{equation}\label{ng:C20}
\frac{z_v-Q}{z_v-Q^{-1}}
                         =\frac{Q+X}{1+QX}.
\end{equation}

This identity follows by factoring the numerator as
\(-(1+Q)(1+X/Q)/(1-X)\) and the denominator as
\(-(1+Q^{-1})(1+XQ)/(1-X)\).

For \(T\), note \(z_w=z_v(P,X^{-1})\). Its ratio on the \(-P\) eigenline of \(BC\) is therefore the inverse of \((P+X)/(1+PX)\). The selected eigenline of \(\Lambda_w=(BC)^{-1}\) with eigenvalue \(-P\) is instead the \(-P^{-1}\) eigenline of \(BC\). Thus its marked ratio is

\begin{equation}\label{ng:C21}
\frac{z_w-P^{-1}}{z_w-P}
                         =\frac{P+X}{1+PX}.
\end{equation}

Equations \eqref{ng:C20}--\eqref{ng:C21} are exactly the two source auxiliary ratios. Their extension to \(Q=P=1\) is the full matrix limit of \eqref{ng:C16}, not a separate evaluation of two coalescing eigenvectors.

The determinants are nonzero nearby. More explicitly, after \eqref{ng:C18},

\begin{equation}\label{ng:C22}
\begin{aligned}
 \det V&=\frac{(\beta+2)(\beta X+X^2+1)}{(X-1)^2},\\
 \det T&=
 -\frac{(X-1)^2(X-M^2)(XM^2-1)
                (\beta M^2+M^4+1)}
 {M^4(\beta X+X^2+1)^2}.
 \end{aligned}
\end{equation}

Both equal \(-4\) at the complete point. These formulas fix the regularity needed for the inverse transport matrices.

\subsection{The main commutator gives the source longitude factor}\label{ng:the-main-commutator-gives-the-source-longitude-factor}

Consider the ordered commutator

\begin{equation}\label{ng:C23}
\Lambda_K=T V^{-1}T^{-1}V.
\end{equation}

It commutes with \(A\). This also follows from the transport paths: its inverse \(V^{-1}TVT^{-1}\) sends \(A\) successively to \(B^{-1}\), \(C\), \(D_0^{-1}\), and back to \(A\), using \eqref{ng:C17}.

Set

\begin{equation}\label{ng:C24}
u=\frac{M^2-X}{1-M^2X}.
\end{equation}

The exact matrix is

\begin{equation}\label{ng:C25}
\Lambda_K=
 \begin{pmatrix}
 -u & (u^{-1}-u)/(M-M^{-1})\\
 0  & -u^{-1}
 \end{pmatrix}.
\end{equation}

The apparent quotient in its upper-right entry is removable at \(M=1\). A direct certificate, without multiplying four complicated inverse matrices, is

\begin{equation}\label{ng:C26}
T V=V
 \begin{pmatrix}
 -u^{-1}&(u-u^{-1})/(M-M^{-1})\\
 0&-u
 \end{pmatrix}T.
\end{equation}

Substitution of \eqref{ng:C7}, \eqref{ng:C16}, and \eqref{ng:C18} makes every entry of the difference zero. Multiplication by \(V^{-1}\) and \(T^{-1}\), which are legitimate by \eqref{ng:C22}, gives \eqref{ng:C25}. The expression for the upper-right entry is regular because \(u=1\) at \(M=1\).

The \(M\) eigenline of the positively oriented puncture \(A\) is the first coordinate line. Thus the eigenvalue of \eqref{ng:C23} on that line is \(-u\), exactly the source factor for one edge. The original source cut selects the opposite \(M^{-1}\) eigenline; its eigenvalue is \(-u^{-1}\), retaining the required inverse-flag convention.

Write \(X=e^{2\pi i x}\), using the logarithm selected by \(x=1/4\)
at the complete edge. At \(M=1\), with \(X\) kept on its continued branch,

\begin{equation}\label{ng:C27}
\Lambda_K=
 \begin{pmatrix}
 -1&-2i\cot(\pi x)\\0&-1
 \end{pmatrix}.
\end{equation}

Since \(A\) translates by one, \eqref{ng:C27} has projective translation
\(2i\cot(\pi x)\). At the complete edge \(x=1/4\), this is \(2i\).
The negative determinant-one lift in \eqref{ng:C25} is the actual commutator lift; it was not inserted by multiplying the word by a sign. No global preferred-longitude formula follows from this local calculation alone.

\subsection{Comparison with the explicit native pure block}\label{ng:comparison-with-the-explicit-native-pure-block}

The face construction in Section~\ref{ng:the-local-face-maps} supplies the geometric marking
for these character matrices. The comparison retains the original
white-map order and all six white/shaded returns. Section~\ref{ng:ordered-pages-and-the-half-twist-collar} then
assembles the ordered tree and its odd seams.

Write its four actual white maps as \(W_0,W_1,W_2,W_3\), and its two plus-copy shaded maps as \(S_v,S_w\). The equatorial corners are \(0,\infty,1,R\). Put

\begin{equation}\label{ng:C28}
N_j=W_{j-1}^{-1}W_j,\qquad j\bmod4.
\end{equation}

The native corner maps have derivatives
\((M^{-2},M^2,M^{-2},M^2)\) in the finite local coordinates, with \(1/z\) at infinity. They obey
\(N_0N_1N_2N_3=I\) by literal cancellation of the white maps.

In the matrix representatives fixed in that construction,

\begin{equation}\label{ng:C29}
\det W_0=L,\quad \det W_1=L/M^2,\quad
 \det W_2=L/M^4,\quad \det W_3=LM^2.
\end{equation}

The determinant-one puncture matrices to compare with \eqref{ng:C6}--\eqref{ng:C13} are therefore

\begin{equation}\label{ng:C30}
A_{\mathrm g}=M^{-3}N_3,\quad
 B_{\mathrm g}=MN_0,\quad
 C_{\mathrm g}=MN_1,\quad
 D_{\mathrm g}=MN_2.
\end{equation}

Their product is one, since the four displayed scalar powers multiply to one. Each trace is \(M+M^{-1}\): its eigenvalue ratio is prescribed by the native corner derivative, its determinant is one, and its trace tends to \(2\) at completeness. Thus the sign is fixed by continuity, rather than reset on a puncture independently.

The two disk words now agree with their actual geometric words, including their determinant-one lifts:

\begin{equation}\label{ng:C31}
A_{\mathrm g}B_{\mathrm g}=M^{-2}W_2^{-1}W_0,\qquad
 (B_{\mathrm g}C_{\mathrm g})^{-1}
                         =M^{-2}W_1^{-1}W_3.
\end{equation}

These are the native \(v,w\) longitudes. Their selected eigenvalues are \(-Q,-P\), with their actual apex flags. The inverse in the second abstract word is essential to \eqref{ng:C31}.

\subsubsection{Actual auxiliary transports and character normalization}\label{ng:actual-auxiliary-transports-and-character-normalization}

The native shaded paths send the corners as follows:

\begin{equation}\label{ng:C32}
S_v:\infty\mapsto0,\quad 1\mapsto R;\qquad
 S_w:R\mapsto0,\quad 1\mapsto\infty.
\end{equation}

Their six white/shaded return identities imply, for the native punctures, that these paths reverse the oriented small boundary meridian at each endpoint. Thus

\begin{equation}\label{ng:C33}
\begin{aligned}
 S_vC_{\mathrm g}S_v^{-1}&=B_{\mathrm g}^{-1},
 &S_vD_{\mathrm g}S_v^{-1}&=A_{\mathrm g}^{-1},\\
 S_wA_{\mathrm g}S_w^{-1}&=B_{\mathrm g}^{-1},
 &S_wD_{\mathrm g}S_w^{-1}&=C_{\mathrm g}^{-1}.
 \end{aligned}
\end{equation}

For example, the projective relation \(W_3S_v\sim S_v^-W_1\), together with
\(W_0S_v\sim S_v^-W_0\), gives
\[
 S_v(W_0^{-1}W_1)S_v^{-1}\sim W_0^{-1}W_3=N_0^{-1}.
\]
The native matrix representatives can differ by a nonconstant scalar in this calculation; those scalars are not set to one. Normalize the punctures using \eqref{ng:C30}. The two sides now have determinant one and trace \(t\), so their projective equality has scalar \(+1\), since \(t\ne0\) near completeness. The remaining three identities follow from the other returns in the same way. This proves the literal determinant-one identities \eqref{ng:C33} without discarding a projective scalar.

Consequently

\begin{equation}\label{ng:C34}
V_{\mathrm g}=S_v^{-1},\qquad
                       T_{\mathrm g}=S_w^{-1}
\end{equation}

satisfy exactly \eqref{ng:C17}. These are the actual positively marked auxiliary meridians of the native construction, not extra operators inserted into the tangle.

There is a regular marked conjugacy to the character coordinates. At completeness the selected \(M\)-eigenlines of \(A_{\mathrm g},B_{\mathrm g}\) are distinct. Their eigenlines vary holomorphically: the matrices minus their specified eigenvalue have rank one with a nonzero entry near the parabolic point. Send those two eigenlines to infinity and zero, then fix the residual diagonal conjugation by setting the upper-right entry of \(A\) to one. This entry is a unit near completeness. The pair becomes exactly \eqref{ng:C6}, since \(\operatorname{tr}(AB)=-\beta\) fixes \(b=-\beta-2\).

The remaining matrix \(C\) is then determined by its trace, the two mixed traces, and \(\operatorname{tr}(ABC)=t\), giving \eqref{ng:C7}. The common centralizer of this irreducible pair consists only of scalars. Hence a transport satisfying \eqref{ng:C17} is unique projectively: the ratio of two such transports commutes with both punctures. Its marked eigenvalue ratio is therefore \eqref{ng:C20}, not a new free parameter.

The branch is fixed by the literal complete matrices. Substituting the complete native white maps gives
\(A_{\mathrm g}=\left(\begin{smallmatrix}3-2i&-1+i\\4-4i&-1+2i\end{smallmatrix}\right)\)
and
\(C_{\mathrm g}=\left(\begin{smallmatrix}1&-1+i\\0&1\end{smallmatrix}\right)\).
Their product has trace \(2+8i\). Since \(S=8\), the character coordinate \((\operatorname{tr}(A_{\mathrm g}C_{\mathrm g})-2)/S\) is \(i\), not the inverse root.

Let \(X_{\mathrm g}\) denote the native source wire extracted from its actual auxiliary meridian:

\begin{equation}\label{ng:C35}
U_v^{\exp}=\frac{Q+X_{\mathrm g}}{1+QX_{\mathrm g}}.
\end{equation}

The character coordinate \(X\) has the same ratio by \eqref{ng:C20}. For \(Q^2\ne1\), the fractional-linear map in \eqref{ng:C35} is injective, so \(X=X_{\mathrm g}\). Holomorphic continuation establishes equality through \(Q=1\). Equivalently,

\begin{equation}\label{ng:C36}
\operatorname{tr}(A_{\mathrm g}C_{\mathrm g})
       =2+\bigl[\beta+\gamma+(M+M^{-1})^2\bigr]X_{\mathrm g}.
\end{equation}

This proves the matching branch, rather than choosing the inverse root because the two have the same discriminant.

\subsubsection{Main basepoint, selected flag, and negative lift}\label{ng:main-basepoint-selected-flag-and-negative-lift}

The native geometric longitude based at infinity is

\begin{equation}\label{ng:C37}
\Lambda_\infty=S_wS_v^{-1}S_w^{-1}S_v.
\end{equation}

Its positively oriented main meridian is \(C_{\mathrm g}=MN_1\). In the native chart it is exactly

\begin{equation}\label{ng:C38}
C_{\mathrm g}=
 \begin{pmatrix}
 M^{-1}&M E_0-M^{-1}\\0&M
 \end{pmatrix},
\end{equation}

where \(E_0\) is the native white-face coordinate. Thus infinity is its actual \(M^{-1}\) source flag. The white-face coordinate \(E_0\) need not equal the source coordinate \(X=X_{\mathrm g}\) away from \(M=1\).

Our abstract commutator is based at the corner \(R\), whose positive meridian is \(A_{\mathrm g}\). The actual path

\begin{equation}\label{ng:C39}
U=S_w^{-1}S_v,\qquad
 \infty\longrightarrow 0\longrightarrow R.
\end{equation}

transports both the meridian and longitude:

\begin{equation}\label{ng:C40}
UC_{\mathrm g}U^{-1}=A_{\mathrm g},\qquad
 U\Lambda_\infty U^{-1}
      =S_w^{-1}S_vS_wS_v^{-1}
      =T_{\mathrm g}V_{\mathrm g}^{-1}
                         T_{\mathrm g}^{-1}V_{\mathrm g}.
\end{equation}

The first identity follows from \eqref{ng:C33}, including its two inversions; the second is literal multiplication of the four words. The source \(M^{-1}\) eigenflag is transported by this same path. It is not identified with infinity in the normalized matrix \eqref{ng:C6}, where infinity is instead the \(M\) eigenline.

Equations \eqref{ng:C25}, \eqref{ng:C36}, and \eqref{ng:C40} identify the native longitude eigenvalues exactly: on its source \(M^{-1}\) eigenflag the value is \(-u^{-1}\); on the preferred \(M\) flag it is \(-u\). The negative sign survives the basepoint change. At completeness the native matrix is

\begin{equation}\label{ng:C41}
\Lambda_\infty=
 \begin{pmatrix}-1&2+2i\\0&-1\end{pmatrix},\qquad
 C_{\mathrm g}=
 \begin{pmatrix}1&i-1\\0&1\end{pmatrix}.
\end{equation}

Their projective translations have ratio \(2i\), equal to \eqref{ng:C27}. This checks the actual marked complete lift and scale; it is not the reason for the full rational identification above.

Thus the two constructions agree on the canonical pure block, its main eigenvalue, its two primitive auxiliary words, and the lift of this local main word. This proves Theorem~\ref{ng:character}. The next section assembles the ordered tree. Identifying a globally preferred main longitude is a separate framing question.
